%%
%% DIRA: Lightweight Machine Learning for Post-Caller Indel Refinement
%% ACM-BCB 2026 Submission
%%
\documentclass[sigconf]{acmart}

\AtBeginDocument{%
  \providecommand\BibTeX{{Bib\TeX}}}

\setcopyright{acmlicensed}
\copyrightyear{2026}
\acmYear{2026}
\acmDOI{10.1145/nnnnnnn.nnnnnnn}
\acmConference[ACM-BCB '26]{15th ACM Conference on Bioinformatics, Computational Biology, and Health Informatics}{2026}{TBD}
\acmISBN{978-1-4503-XXXX-X/2026/XX}

\usepackage{booktabs}
\usepackage{multirow}
\usepackage{amsmath}

\begin{document}

\title{Lightweight Discriminative Indel Refinement via Artifact-Aware Alignment Modeling}

\author{Author A}
\authornote{All authors contributed equally to this research.}
\affiliation{%
  \institution{Institution}
  \country{Country}}
\email{authora@institution.edu}

\author{Author B}
\authornotemark[1]
\affiliation{%
  \institution{Institution}
  \country{Country}}
\email{authorb@institution.edu}

\author{Author C}
\authornotemark[1]
\affiliation{%
  \institution{Institution}
  \country{Country}}
\email{authorc@institution.edu}

\renewcommand{\shortauthors}{Author A et al.}

\begin{abstract}
Accurate detection of insertions and deletions (indels) from short-read sequencing data remains challenging due to alignment artifacts, homopolymer errors, and systematic sequencing biases. While assembly-based callers and deep learning methods achieve high accuracy, they incur substantial computational overhead. We present DIRA, a lightweight post-calling refinement framework that improves indel accuracy from conventional variant callers without modifying the caller itself. DIRA extracts 43 alignment-based features from BAM files at each candidate indel site, including novel haplotype consistency features that approximate assembly-level signals through CIGAR pattern analysis, and uses a gradient-boosted decision tree classifier to distinguish true indels from artifacts. A key methodological contribution is CIGAR-aware read classification, which correctly partitions reads into variant-supporting and reference groups by examining insertion/deletion operations in read alignments rather than relying on base-level matching. On the Genome in a Bottle HG002 benchmark, DIRA improves bcftools indel F1 from 0.8100 to 0.8308 on a fully held-out chromosome, with a 19.4\% reduction in false positives. When applied to FreeBayes candidates---without retraining---DIRA achieves an 8.4\% absolute F1 improvement (0.7511 to 0.8347), demonstrating caller-agnostic artifact detection. Cross-sample evaluation on HG003 confirms generalization. The pipeline adds under 25 minutes of compute for feature extraction and under 5 seconds for prediction on an 8-core cloud instance, making it suitable as a drop-in refinement layer for production genomics pipelines.
\end{abstract}

\begin{CCSXML}
<ccs2012>
 <concept>
  <concept_id>10010405.10010444.10010450</concept_id>
  <concept_desc>Applied computing~Computational genomics</concept_desc>
  <concept_significance>500</concept_significance>
 </concept>
 <concept>
  <concept_id>10010405.10010444.10010448</concept_id>
  <concept_desc>Applied computing~Bioinformatics</concept_desc>
  <concept_significance>500</concept_significance>
 </concept>
 <concept>
  <concept_id>10010147.10010257.10010293.10010294</concept_id>
  <concept_desc>Computing methodologies~Supervised learning by classification</concept_desc>
  <concept_significance>300</concept_significance>
 </concept>
</ccs2012>
\end{CCSXML}

\ccsdesc[500]{Applied computing~Computational genomics}
\ccsdesc[500]{Applied computing~Bioinformatics}
\ccsdesc[300]{Computing methodologies~Supervised learning by classification}

\keywords{variant calling, indel detection, machine learning, gradient boosting, sequence alignment, bioinformatics, variant filtering}

\maketitle

%% ========================================================================
\section{Introduction}
%% ========================================================================

Accurate identification of insertions and deletions (indels) from next-generation sequencing (NGS) data is essential for clinical genomics, population genetics, and functional annotation of genomes~\cite{goldfeder2016medical}. While single nucleotide polymorphism (SNP) calling has achieved near-perfect accuracy, indel calling remains substantially more difficult due to the complex interaction between alignment artifacts, sequencing errors, and variant representation ambiguity~\cite{zook2019open, li2011statistical}.

Current approaches to indel calling span a wide spectrum of computational cost and accuracy. At one extreme, pileup-based callers such as bcftools~\cite{li2011statistical} and FreeBayes~\cite{garrison2012haplotype} offer fast variant detection but produce significant numbers of false positives, particularly in homopolymer runs, tandem repeats, and low-mappability regions. At the other extreme, assembly-based methods such as GATK HaplotypeCaller~\cite{poplin2018scaling} perform local de Bruijn graph assembly and PairHMM likelihood computation, achieving substantially higher accuracy at the cost of increased runtime. More recently, deep learning approaches such as DeepVariant~\cite{poplin2018universal} have demonstrated that convolutional neural networks can learn complex error signatures directly from pileup images, while hybrid approaches such as DNAscope~\cite{freed2022dnascope} combine assembly mathematics with gradient-boosted machine learning for genotyping.

A notable gap exists between these categories: lightweight callers remain widely used in large-scale studies and clinical settings where computational resources are constrained, yet their indel accuracy lags significantly behind assembly and deep learning methods. Post-calling variant quality score recalibration (VQSR) offers limited improvement and has been shown to sometimes decrease indel accuracy~\cite{zhao2020accuracy}.

We present DIRA (Deep Indel Refinement via Alignment features), a post-calling framework that addresses this gap. DIRA takes candidate indels from any upstream caller, extracts alignment-based features directly from BAM files, and uses a LightGBM~\cite{ke2017lightgbm} classifier to identify and filter likely artifacts. The key insight is that alignment artifacts exhibit characteristic signatures---inconsistent CIGAR patterns across reads, elevated edit distances, strand bias, fragment discordance---that can be captured by carefully engineered features without performing local assembly.

Our main contributions are:
\begin{enumerate}
    \item \textbf{CIGAR-aware read classification}, which correctly partitions reads into indel-supporting and reference groups by examining CIGAR operations rather than base-level matching, enabling accurate computation of all per-group features.
    \item \textbf{Haplotype consistency features} that approximate assembly-level signals by measuring CIGAR pattern agreement and entropy across reads at each site, inspired by DNAscope's assembly entropy annotation~\cite{freed2022dnascope}.
    \item \textbf{Caller-agnostic artifact detection}: a model trained on bcftools candidates that transfers to FreeBayes without retraining, achieving an 8.4\% absolute F1 improvement.
    \item \textbf{A modular, efficient pipeline} that separates feature extraction from inference, completing whole-genome feature extraction in under 25 minutes and prediction in under 5 seconds on a standard cloud instance.
\end{enumerate}

%% ========================================================================
\section{Related Work}
%% ========================================================================

\subsection{Pileup-Based Variant Callers}
bcftools~\cite{li2011statistical} uses a Bayesian framework with allele frequencies estimated from sequencing data, while FreeBayes~\cite{garrison2012haplotype} employs a haplotype-based Bayesian approach without local assembly. Both are computationally efficient but rely on simplified statistical models that do not capture complex sequencing error modes, leading to elevated false positive rates for indels.

\subsection{Assembly-Based and Hybrid Callers}
GATK HaplotypeCaller~\cite{poplin2018scaling} identifies active regions, performs local de Bruijn graph assembly to enumerate candidate haplotypes, and uses PairHMM to compute read-haplotype likelihoods. DNAscope~\cite{freed2022dnascope} extends this architecture by replacing the Bayesian genotyping model with a gradient-boosted machine learning model trained on variant annotations including assembly entropy, neighbourhood base quality (NBQ), and read position distance. DNAscope achieves 99.46\% indel F1 on HG002 at 30x coverage. Strelka2~\cite{kim2018strelka2} uses a tiered haplotype model with candidate indel assembly for germline and somatic calling.

\subsection{Deep Learning Approaches}
DeepVariant~\cite{poplin2018universal} encodes the reference and aligned reads at each candidate site as a multi-channel pileup image and uses an Inception-architecture CNN~\cite{szegedy2015rethinking} to predict diploid genotype likelihoods. The approach achieves state-of-the-art accuracy but requires substantial compute resources and GPU inference. Several subsequent approaches have extended this paradigm~\cite{luo2018clairvoyante}.

\subsection{Variant Filtering and Recalibration}
GATK's Variant Quality Score Recalibration (VQSR) applies a Gaussian mixture model to variant annotations to identify likely true variants. However, VQSR has been reported to decrease indel F1 in some settings~\cite{zhao2020accuracy}. Hard filtering based on quality thresholds remains common but lacks the expressiveness to capture complex artifact signatures.

\subsection{Position of DIRA}
DIRA is most closely related to DNAscope's approach---using gradient-boosted trees on engineered features---but operates \textit{post-calling} rather than within the variant caller. This means DIRA does not have access to assembly-level information (de Bruijn graphs, PairHMM likelihoods). To compensate, DIRA approximates assembly-level signals through CIGAR pattern analysis and read consistency metrics computed directly from BAM alignments.

%% ========================================================================
\section{Methods}
%% ========================================================================

\subsection{Overview}
DIRA operates as a four-stage pipeline: (1) variant calling with a conventional caller, (2) feature extraction from BAM alignments at candidate indel sites, (3) model training with truth set labels, and (4) inference and VCF annotation. Stages are implemented as independent scripts, allowing modular execution and separate optimization.

\subsection{CIGAR-Aware Read Classification}
\label{sec:cigar}
A fundamental challenge in computing per-group features (e.g., mapping quality of variant-supporting reads vs.\ reference reads) at indel sites is correctly classifying which reads support the variant allele. A naive approach examines the base at the pileup position and classifies reads based on whether the base matches the ALT allele. This heuristic is unreliable for indels: reads carrying an insertion encode it as a CIGAR \texttt{I} operation, and the base at the pileup column may match the reference. Similarly, reads spanning a deletion show a CIGAR \texttt{D} operation that is invisible to base-level inspection.

DIRA instead walks each read's CIGAR string and classifies a read as indel-supporting if it contains an insertion (\texttt{I}) or deletion (\texttt{D}) operation within $\pm 2$\,bp of the variant position (to account for left-alignment differences), with length validation ensuring the CIGAR operation length approximately matches the expected indel size:

\begin{equation}
|l_{\text{CIGAR}} - |l_{\text{indel}}|| \leq \max\left(1, \left\lfloor\frac{|l_{\text{indel}}|}{2}\right\rfloor\right)
\end{equation}

This classification is applied to all reads overlapping the variant site, collected via \texttt{bam.fetch()} rather than \texttt{bam.pileup()} to avoid the overhead of full pileup column construction.

\subsection{Feature Engineering}
DIRA extracts 43 features per candidate indel, organized into seven categories. All per-group features (mapping quality, base quality, strand ratio, etc.) use the CIGAR-aware classification described in Section~\ref{sec:cigar}.

\subsubsection{Depth and Allele Balance (3 features).}
Total depth, alternate allele depth, and allele balance (ratio of alt to total reads). These features capture the basic evidence for the variant.

\subsubsection{Mapping Quality (5 features).}
Mean and variance of mapping quality (MQ) for alt and reference reads, fraction of reads with MQ\,=\,0, and the MQ difference between alt and ref groups. Artifact-supporting reads often have degraded mapping quality~\cite{li2008mapping}. The MQ\,=\,0 fraction captures unmappable reads indicative of repetitive regions.

\subsubsection{Strand and Position Bias (6 features).}
Forward strand ratio for alt and ref reads, mean and variance of read-end distance (minimum of distance to either read end) for all reads, and stratified read-end distance for alt and ref groups separately. Alignment artifacts frequently cluster at read ends~\cite{freed2022dnascope} and exhibit strand asymmetry.

\subsubsection{Base Quality (8 features).}
Mean and variance of whole-read base quality for alt and ref reads, neighbourhood base quality (NBQ)---mean BQ within $\pm 5$\,bp of the variant position~\cite{challis2012atlas2}---for alt and ref reads, the BQ drop between ref and alt groups, and flanking base quality in $\pm 5$\,bp computed separately for alt and ref reads with actual reference comparison for mismatch detection.

\subsubsection{Alignment Quality (6 features).}
Mean NM (edit distance) tag~\cite{li2009sequence} for alt and ref reads, proper pair rate for alt and ref reads, and mismatch rate near the variant site computed via direct reference comparison. Edit distance is the second most important feature in our model, consistent with the observation that artifact-supporting reads exhibit globally elevated mismatch rates.

\subsubsection{Haplotype Consistency (3 features).}
These features, inspired by DNAscope's assembly entropy annotation~\cite{freed2022dnascope}, approximate assembly-level signals without de Bruijn graph construction:
\begin{itemize}
    \item \textbf{CIGAR agreement}: for each alt read, we extract a CIGAR signature---a tuple of (operation, length) pairs within $\pm 3$\,bp of the variant position. The CIGAR agreement score is the fraction of alt reads sharing the most common signature. True indels produce consistent signatures ($>0.8$); artifacts produce diverse patterns.
    \item \textbf{Haplotype entropy}: Shannon entropy~\cite{shannon1948mathematical} of CIGAR signatures across alt reads. This directly approximates assembly graph entropy---low entropy indicates convergence to a single haplotype, while high entropy indicates conflicting read structures.
    \item \textbf{Fragment support rate}: for paired-end data, the fraction of alt reads whose mate is also classified as indel-supporting. True genomic variants produce concordant fragment-level evidence; artifacts often appear on only one mate of a fragment.
\end{itemize}

\subsubsection{Sequence Context and Indel Properties (12 features).}
Soft-clip rates for alt and ref reads and their difference, indel length and direction, insertion sequence entropy, homopolymer run length at the variant position, GC content in a 100\,bp window, tandem repeat unit length and total tract length (capturing STR context beyond homopolymers), insert size statistics for alt and ref reads, and the UDP/DP ratio (unfiltered to filtered depth)~\cite{freed2022dnascope}.

\subsection{Performance Optimization}
Feature extraction---the pipeline's computational bottleneck---is parallelized across chromosomes using Python multiprocessing, with one worker per chromosome. Each worker opens independent BAM and FASTA handles (pysam and cyvcf2 objects are not fork-safe). Additional optimizations include: pure-Python statistics to avoid NumPy overhead on small arrays ($n < 100$), early-exit NBQ computation exploiting the sorted order of aligned pairs, set-based $O(1)$ read classification lookup, and batched CSV writes. Feature extraction for 22 autosomes ($\sim$1M indels) completes in 30--40 minutes on a GCP n2-standard-8 instance (8 vCPU, 32\,GB RAM).

Prediction is fully decoupled from extraction: the trained model scores a pre-extracted CSV in a single batched call, completing in under 5 seconds regardless of dataset size.

\subsection{Model Training}
We use LightGBM~\cite{ke2017lightgbm}, a gradient-boosted decision tree framework, with binary cross-entropy loss. Key hyperparameters: learning rate 0.05, 64 leaves, feature fraction 0.8, bagging fraction 0.8 with frequency 5, minimum 50 samples per leaf. Training uses early stopping with patience 50 on validation loss.

Class imbalance is handled via \texttt{scale\_pos\_weight} set to the negative/positive ratio. Platt calibration~\cite{platt1999probabilistic} is fitted on the validation set's raw scores to produce calibrated probabilities used for Phred-scaled QUAL scores: $\text{QUAL} = -10 \log_{10}(1 - P)$. The decision threshold $\tau$ is optimized by maximizing F1 on the precision-recall curve.

\subsection{Data Splitting}
To prevent spatial leakage from nearby variants sharing local sequence context, we split by chromosome rather than randomly: chromosomes 1--14 for training, 15--16 for validation, and 17--18 for test. Chromosome 20 is held out entirely for final benchmarking.

\subsection{Experimental Setup}
\begin{itemize}
    \item \textbf{Reference genome}: GRCh38
    \item \textbf{Training sample}: GIAB HG002, 2$\times$250\,bp Illumina, $\sim$30x coverage
    \item \textbf{Truth set}: GIAB v4.2.1 high-confidence indel calls~\cite{zook2019open, wagner2022benchmarking}
    \item \textbf{Evaluation}: RTG Tools vcfeval~\cite{krusche2019best}
    \item \textbf{Candidate generation}: bcftools mpileup/call and FreeBayes
    \item \textbf{Cross-sample evaluation}: GIAB HG003
\end{itemize}

%% ========================================================================
\section{Results}
%% ========================================================================

\subsection{Development Benchmark (HG002, Chromosome 20)}
All development results use RTG vcfeval on chromosome 20, which was fully held out from training, validation, and threshold selection. Table~\ref{tab:dev} shows the progression of DIRA accuracy across feature set iterations.

\begin{table}[h]
\caption{Development results: bcftools indels on HG002 chr20, evaluated with RTG vcfeval. Each row represents a feature extractor iteration.}
\label{tab:dev}
\begin{tabular}{lcccc}
\toprule
\textbf{Method} & \textbf{Precision} & \textbf{Recall} & \textbf{F1} & \textbf{FP} \\
\midrule
bcftools (raw) & 0.7962 & 0.8243 & 0.8100 & 2{,}856 \\
DIRA v4 (base ML) & 0.8011 & 0.8294 & 0.8150 & 2{,}788 \\
DIRA v5 (CIGAR) & 0.8184 & 0.8336 & 0.8260 & 2{,}503 \\
DIRA v6 (full) & \textbf{0.8302} & \textbf{0.8314} & \textbf{0.8308} & \textbf{2{,}302} \\
\bottomrule
\end{tabular}
\end{table}

DIRA v6 achieves a 2.08\% absolute F1 improvement over raw bcftools, with precision increasing by 3.4\% and a 19.4\% reduction in false positives (2{,}856 $\to$ 2{,}302). Importantly, recall is maintained (0.8243 $\to$ 0.8314), indicating that the model identifies artifacts without discarding true variants.

\subsection{Impact of CIGAR-Aware Classification}
The transition from base-level matching (v4) to CIGAR-aware classification (v5) produced the largest single improvement: F1 increased from 0.8150 to 0.8260 (+1.1\%). This validates the hypothesis that correct read partitioning is essential for informative per-group features.

\subsection{Cross-Caller Generalization (FreeBayes)}
To test whether DIRA learns caller-agnostic artifact signatures, we applied the model---trained exclusively on bcftools candidates from HG002---to FreeBayes whole-genome output from HG002 without retraining.

\begin{table}[h]
\caption{Cross-caller evaluation: FreeBayes indels on HG002 whole genome. The DIRA model was trained on bcftools candidates only.}
\label{tab:freebayes}
\begin{tabular}{lccc}
\toprule
\textbf{Method} & \textbf{Precision} & \textbf{Recall} & \textbf{F1} \\
\midrule
FreeBayes (raw) & 0.6550 & 0.8801 & 0.7511 \\
FreeBayes + DIRA & \textbf{0.8306} & 0.8388 & \textbf{0.8347} \\
\midrule
$\Delta$ & +17.6\% & $-$4.1\% & \textbf{+8.4\%} \\
\bottomrule
\end{tabular}
\end{table}

DIRA removes 179{,}944 false positives (63.2\% FP reduction) while losing only 4.1\% recall. The 8.4\% absolute F1 gain demonstrates that alignment artifact signatures are shared across callers, and that DIRA's features capture caller-independent error modes.

\subsection{Cross-Sample Generalization (HG003)}

%% PLACEHOLDER: Insert HG003 results when available
%% Expected format:
%% \begin{table}[h]
%% \caption{Cross-sample evaluation: HG003, FreeBayes + DIRA.}
%% \label{tab:hg003}
%% \begin{tabular}{lccc}
%% \toprule
%% Method & Precision & Recall & F1 \\
%% \midrule
%% FreeBayes (raw) & X.XXXX & X.XXXX & X.XXXX \\
%% FreeBayes + DIRA & X.XXXX & X.XXXX & X.XXXX \\
%% bcftools (raw) & X.XXXX & X.XXXX & X.XXXX \\
%% bcftools + DIRA & X.XXXX & X.XXXX & X.XXXX \\
%% \bottomrule
%% \end{tabular}
%% \end{table}

\textit{[HG003 cross-sample results to be inserted upon completion of evaluation.]}

\subsection{Feature Importance}
Table~\ref{tab:importance} shows the top 10 features by LightGBM gain importance. The NM (edit distance) features are the strongest non-depth signals, consistent with the observation that artifact-supporting reads exhibit elevated mismatch rates. The CIGAR agreement feature---introduced in v5---ranks second, validating the haplotype consistency approach. Allele balance, which was uninformative under base-level matching (v4), becomes the third most important feature under CIGAR-aware classification, confirming that correct read partitioning is necessary for this feature to carry signal.

\begin{table}[h]
\caption{Top 10 features by LightGBM gain importance.}
\label{tab:importance}
\begin{tabular}{clc}
\toprule
\textbf{Rank} & \textbf{Feature} & \textbf{Category} \\
\midrule
1 & nm\_mean\_alt & Alignment quality \\
2 & cigar\_agreement & Haplotype consistency \\
3 & allele\_balance & Depth \\
4 & is\_insertion & Indel properties \\
5 & mismatch\_rate\_alt & Alignment quality \\
6 & indel\_length & Indel properties \\
7 & mq\_mean\_alt & Mapping quality \\
8 & nm\_mean\_ref & Alignment quality \\
9 & bq\_mean\_alt & Base quality \\
10 & softclip\_rate\_alt & Soft clipping \\
\bottomrule
\end{tabular}
\end{table}

\subsection{Training Metrics}
With chromosome-based splitting (train: chr1--14, val: chr15--16, test: chr17--18), the v6 model achieves validation F1 = 0.8828 (AUROC = 0.8982) and test F1 = 0.8804 (AUROC = 0.8875). The close agreement between validation and test metrics indicates that the chromosome-based split produces consistent generalization behavior.

\subsection{Runtime Analysis}
Table~\ref{tab:runtime} summarizes DIRA's computational requirements on a GCP n2-standard-8 instance (8 vCPU, 32\,GB RAM, Intel Xeon).

\begin{table}[h]
\caption{Runtime for whole-genome HG002 (chr1--22, $\sim$1M indels).}
\label{tab:runtime}
\begin{tabular}{lcc}
\toprule
\textbf{Stage} & \textbf{Time} & \textbf{Notes} \\
\midrule
Feature extraction & 30--40 min & 8 workers, parallelized \\
Model training & $\sim$2 min & 832 boosting rounds \\
Prediction & $<$5 sec & Batched, no BAM access \\
\midrule
Total overhead & $\sim$35 min & On top of caller runtime \\
\bottomrule
\end{tabular}
\end{table}

%% ========================================================================
\section{Discussion}
%% ========================================================================

\subsection{Comparison to State of the Art}
DIRA's absolute indel F1 (0.8308 on chr20) is substantially below assembly-based and deep learning callers: DNAscope achieves 99.46\%~\cite{freed2022dnascope} and DeepVariant achieves 99.51\%~\cite{poplin2018universal} on similar benchmarks. However, this comparison is not direct: DIRA is a post-calling filter bounded by the upstream caller's candidate sensitivity, while DNAscope and DeepVariant perform their own candidate generation with local assembly. The relevant comparison is DIRA-filtered bcftools versus raw bcftools, where DIRA provides meaningful improvement without modifying the caller.

The cross-caller result is perhaps more significant: DIRA achieves a larger absolute gain on FreeBayes (+8.4\% F1) than on bcftools (+2.1\% F1), consistent with the expectation that noisier callers benefit more from post-calling refinement. This suggests DIRA is most valuable as a lightweight cleanup layer for high-sensitivity, low-specificity callers.

\subsection{Why CIGAR-Aware Classification Matters}
At an insertion site, inserted bases do not appear in the pileup column---they exist only in the CIGAR string. Base-level matching therefore produces near-random alt/ref partitioning at many indel sites, contaminating all stratified features with noise. The 1.1\% F1 improvement from CIGAR-aware classification (v4 $\to$ v5) reflects the fundamental importance of correct read classification for meaningful feature computation.

\subsection{Approximating Assembly Entropy}
DNAscope's assembly entropy measures disorder in the de Bruijn graph at active regions~\cite{freed2022dnascope}. Our CIGAR signature entropy approximates this by measuring disorder in read-level evidence. While not equivalent---CIGAR signatures are noisier than assembled haplotypes---the feature ranks second by importance, confirming it captures useful signal about local alignment consistency without the computational cost of graph construction.

\subsection{Limitations}
\textbf{Bounded by caller sensitivity.} DIRA cannot recover indels that the upstream caller fails to detect. The recall ceiling is the caller's raw recall, which limits absolute F1 regardless of model quality.

\textbf{Single-sample training.} The model is trained on HG002 only. While cross-caller generalization is strong, cross-sample and cross-platform generalization requires further evaluation.

\textbf{No local assembly.} Without assembly, DIRA cannot resolve complex multi-allelic indels, overlapping events, or sites where the variant representation is fundamentally ambiguous.

\textbf{Simplified mismatch detection.} The mismatch rate feature compares read bases to the reference without accounting for known germline SNPs, potentially inflating counts at heterozygous sites.

%% ========================================================================
\section{Conclusion}
%% ========================================================================

We presented DIRA, a lightweight post-calling framework for indel refinement that achieves meaningful accuracy improvements over conventional variant callers. The key technical insights are: (1) CIGAR-aware read classification is essential for computing informative per-group features at indel sites, (2) assembly-level signals can be approximated through CIGAR pattern analysis and fragment concordance without performing local assembly, and (3) alignment artifact signatures are shared across variant callers, enabling cross-caller transfer without retraining.

DIRA is designed as a drop-in refinement layer for existing pipelines: it requires no modification to the upstream caller, adds minimal compute overhead, and is compatible with standard bioinformatics tools. Future work will focus on multi-sample training for improved generalization, integration with additional callers, and lightweight micro-assembly at candidate sites to generate true haplotype features.

%% ========================================================================
\begin{acks}
This work was performed using Google Cloud Platform compute resources. We acknowledge the Genome in a Bottle Consortium and the National Institute of Standards and Technology for providing the benchmark datasets used in this study.
\end{acks}

\bibliographystyle{ACM-Reference-Format}
\begin{thebibliography}{20}

\bibitem{li2011statistical}
Heng Li.
\newblock A statistical framework for {SNP} calling, mutation discovery, association mapping and population genetical parameter estimation from sequencing data.
\newblock {\em Bioinformatics}, 27(21):2987--2993, 2011.

\bibitem{garrison2012haplotype}
Erik Garrison and Gabor Marth.
\newblock Haplotype-based variant detection from short-read sequencing.
\newblock {\em arXiv preprint arXiv:1207.3907}, 2012.

\bibitem{poplin2018scaling}
Ryan Poplin, Valentin Ruber-Vogelgesang, et~al.
\newblock Scaling accurate genetic variant discovery to tens of thousands of samples.
\newblock {\em bioRxiv}, 2018.

\bibitem{poplin2018universal}
Ryan Poplin, Pi-Chuan Chang, David Alexander, et~al.
\newblock A universal {SNP} and small-indel variant caller using deep neural networks.
\newblock {\em Nature Biotechnology}, 36(10):983--987, 2018.

\bibitem{freed2022dnascope}
Donald Freed, Renke Pan, Haodong Chen, et~al.
\newblock {DNAscope}: High accuracy small variant calling using machine learning.
\newblock {\em bioRxiv}, 2022.
\newblock doi:10.1101/2022.05.20.492556.

\bibitem{ke2017lightgbm}
Guolin Ke, Qi~Meng, Thomas Finley, et~al.
\newblock {LightGBM}: A highly efficient gradient boosting decision tree.
\newblock In {\em Advances in Neural Information Processing Systems}, 2017.

\bibitem{kim2018strelka2}
Sangtae Kim, Konrad Scheffler, Aaron~L Halpern, et~al.
\newblock Strelka2: fast and accurate calling of germline and somatic variants.
\newblock {\em Nature Methods}, 15(8):591--594, 2018.

\bibitem{zook2019open}
Justin~M Zook, Jennifer McDaniel, Nathan~D Olson, et~al.
\newblock An open resource for accurately benchmarking small variant and reference calls.
\newblock {\em Nature Biotechnology}, 37(5):561--566, 2019.

\bibitem{wagner2022benchmarking}
Justin Wagner, Nathan~D Olson, Lindsay Harris, et~al.
\newblock Benchmarking challenging small variants with linked and long reads.
\newblock {\em Cell Genomics}, 2(5), 2022.

\bibitem{krusche2019best}
Peter Krusche, Len Trigg, Paul~C Boutros, et~al.
\newblock Best practices for benchmarking germline small-variant calls in human genomes.
\newblock {\em Nature Biotechnology}, 37(5):555--560, 2019.

\bibitem{goldfeder2016medical}
Rachel~L Goldfeder, Eric~A Priest, Justin~M Zook, et~al.
\newblock Medical implications of technical accuracy in genome sequencing.
\newblock {\em Genome Medicine}, 8(1):24, 2016.

\bibitem{zhao2020accuracy}
Sen Zhao, Oleg Agafonov, Abdulrahman Azab, et~al.
\newblock Accuracy and efficiency of germline variant calling pipelines for human genome data.
\newblock {\em Scientific Reports}, 10:20222, 2020.

\bibitem{challis2012atlas2}
Danny Challis, Jin Yu, Uday~S Evani, et~al.
\newblock An integrative variant analysis suite for whole exome next-generation sequencing data.
\newblock {\em BMC Bioinformatics}, 13:8, 2012.

\bibitem{li2008mapping}
Heng Li, Jue Ruan, and Richard Durbin.
\newblock Mapping short {DNA} sequencing reads and calling variants using mapping quality scores.
\newblock {\em Genome Research}, 18(11):1851--1858, 2008.

\bibitem{li2009sequence}
Heng Li, Bob Handsaker, Alec Wysoker, et~al.
\newblock The {S}equence {A}lignment/{M}ap format and {SAM}tools.
\newblock {\em Bioinformatics}, 25(16):2078--2079, 2009.

\bibitem{platt1999probabilistic}
John Platt.
\newblock Probabilistic outputs for support vector machines and comparisons to regularized likelihood methods.
\newblock {\em Advances in Large Margin Classifiers}, pages 61--74, 1999.

\bibitem{shannon1948mathematical}
Claude~E Shannon.
\newblock A mathematical theory of communication.
\newblock {\em The Bell System Technical Journal}, 27(3):379--423, 1948.

\bibitem{szegedy2015rethinking}
Christian Szegedy, Vincent Vanhoucke, Sergey Ioffe, et~al.
\newblock Rethinking the inception architecture for computer vision.
\newblock In {\em Proceedings of the IEEE Conference on Computer Vision and Pattern Recognition}, pages 2818--2826, 2016.

\bibitem{luo2018clairvoyante}
Ruibang Luo, Fritz~J Sedlazeck, Tak-Wah Lam, and Michael~C Schatz.
\newblock Clairvoyante: a multi-task convolutional deep neural network for variant calling in single molecule sequencing.
\newblock {\em bioRxiv}, 2018.

\bibitem{friedman2001greedy}
Jerome~H Friedman.
\newblock Greedy function approximation: a gradient boosting machine.
\newblock {\em Annals of Statistics}, 29(5):1189--1232, 2001.

\end{thebibliography}

\end{document}
